% source: manti.01.ma, 1ère SM, série suites (Drive 0B6pfSJcEVng0RzBRSWM0UkhSekk)
\ex{Suites}{RzBR}{Ex. 1}{9}
On considère la suite $(U_n)_{n\in\mathbb{N}}$ définie par $U_0=2$ et $U_{n+1}=\sqrt{\dfrac12U_n^2+1}$.
\begin{enumerate}
\item Montrer que $(\forall n\in\mathbb{N})\ U_n\geq\sqrt2$.
\item Étudier la monotonie de la suite $(U_n)_{n\in\mathbb{N}}$.
\item On pose $W_n=U_n^2-2$. Montrer que $(W_n)_n$ est une suite géométrique et en déduire que $U_n=\sqrt2\sqrt{1+\left(\dfrac12\right)^n}$.
\end{enumerate}
\fin
\ex{Suites}{RzBR}{Ex. 2}{9}
On considère la suite $(U_n)_{n\geq1}$ définie par $U_1=\dfrac12$ et $U_{n+1}=\dfrac{nU_n+1}{n+1}$.
\begin{enumerate}
\item Montrer que $(\forall n\in\mathbb{N}^*)\ U_n<1$.
\item On pose $V_n=nU_n$.
\begin{enumerate}
\item Montrer que $(V_n)_{n\geq1}$ est une suite arithmétique et calculer $U_n$ en fonction de $n$.
\item Calculer en fonction de $n$ la somme $S=U_1+2U_2+3U_3+\dots+nU_n$.
\end{enumerate}
\end{enumerate}
\fin
\ex{Suites}{RzBR}{Ex. 3}{8}
Soit $(U_n)_{n\in\mathbb{N}}$ la suite définie par $U_0=2$ et $U_{n+1}=\dfrac14U_n+\dfrac12\left(\dfrac34\right)^n$.
\begin{enumerate}
\item Calculer $U_1$ et $U_2$.
\item On pose $V_n=U_n-\left(\dfrac34\right)^n$.
\begin{enumerate}
\item Montrer que $(V_n)_{n\geq0}$ est une suite géométrique.
\item Calculer $U_n$ en fonction de $n$.
\end{enumerate}
\item Calculer la somme $T=U_0+U_1+U_2+\dots+U_n$.
\end{enumerate}
\fin
\ex{Suites}{RzBR}{Ex. 4}{9}
Soit $(U_n)_{n\in\mathbb{N}}$ la suite définie par $U_0=\dfrac12$ et $U_{n+1}=\dfrac{2U_n}{\sqrt{4+2U_n^2}}$.
\begin{enumerate}
\item Montrer que $(\forall n\in\mathbb{N})\ 0<U_n$.
\item Étudier la monotonie de la suite $(U_n)_{n\in\mathbb{N}}$.
\item On pose $V_n=\dfrac{4}{U_n^2}$. Montrer que $(V_n)_{n\geq0}$ est une suite arithmétique et calculer $V_n$ en fonction de $n$.
\item En déduire que $U_n=\dfrac{2}{\sqrt{2n+16}}$ et calculer $S=\dfrac{1}{U_0^2}+\dfrac{1}{U_1^2}+\dots+\dfrac{1}{U_n^2}$.
\end{enumerate}
\fin
\ex{Suites}{RzBR}{Ex. 5}{10}
Soient $(U_n)_n$ et $(V_n)_n$ les suites définies par
$U_0=1,\ U_{n+1}=\dfrac{U_n+2V_n}{3}$ \quad et \quad $V_0=12,\ V_{n+1}=\dfrac{U_n+3V_n}{4}$.
On pose $W_n=V_n-U_n$ et $T_n=3U_n+8V_n$.
\begin{enumerate}
\item Montrer que $(W_n)_n$ est une suite géométrique et calculer $W_n$ en fonction de $n$.
\item Montrer que $(T_n)_n$ est une suite constante et préciser sa valeur.
\item En déduire $U_n$ et $V_n$ en fonction de $n$.
\end{enumerate}
\fin
\ex{Suites}{RzBR}{Ex. 6}{9}
Soit $(U_n)_{n\in\mathbb{N}}$ la suite définie par $U_0=2$ et $U_{n+1}=\dfrac12U_n+2n$. On pose $V_n=U_n-4n+8$.
\begin{enumerate}
\item Calculer $U_1$ et montrer par récurrence que $(\forall n\in\mathbb{N}^*)\ U_n\geq n$.
\item Montrer que $(V_n)_n$ est une suite géométrique en précisant sa raison.
\item Calculer $U_n$ en fonction de $n$.
\item Calculer $S_n=V_0+V_1+\dots+V_{n-1}$ en fonction de $n$, puis en déduire $T_n=U_0+U_1+\dots+U_{n-1}$ en fonction de $n$.
\end{enumerate}
\fin
\ex{Suites}{RzBR}{Ex. 7}{9}
\begin{enumerate}
\item Soit $(U_n)_{n\in\mathbb{N}}$ une suite arithmétique de raison $r\neq0$ telle que $U_0=2$ et que $U_1$, $U_3$, $U_{13}$ soient des termes consécutifs d'une suite géométrique.
Montrer que $r=-4$ et calculer la somme $U_0+U_1+\dots+U_{n-1}$.
\item Soit $(V_n)_n$ une suite géométrique telle que $V_0V_1V_2=8$ et $V_0+V_1+V_2=7$.
Montrer que $V_1=2$ et en déduire que $q=2$ ou $q=\dfrac12$.
\item Soit $(U_n)_{n\in\mathbb{N}}$ la suite définie par $U_0=6$, $U_1=5$ et $6U_{n+2}=7U_{n+1}-2U_n$. On pose $x_n=U_{n+1}-kU_n$.
Déterminer $k$ pour que $(x_n)_n$ soit une suite géométrique et préciser sa raison, puis déterminer $U_n$ en fonction de $n$.
\end{enumerate}
\fin
